\input{../tex-inputs/latex-leseansicht-vorspann}

\section[Arthur und Olga Schnitzler an Gustav Schwarzkopf, 26. 6. 1908]{L04373 Arthur und Olga Schnitzler an Gustav Schwarzkopf, 26. 6. 1908}
\nopagebreak\mylabel{L04373v}
\rehead{ }\normalsize\beginnumbering\briefempfaengerindex{Schwarzkopf, Gustav@\textsc{Schwarzkopf, Gustav}!zzzSchnitzler, Olga@\emph{von Olga Schnitzler}!1908-06-262@{26. 6. 1908}|(be}\briefempfaengerindex{Schwarzkopf, Gustav@\textsc{Schwarzkopf, Gustav}!zzzSchnitzler, Arthur@\emph{von Arthur Schnitzler}!1908-06-262@{26. 6. 1908}|(be}
\toendnotes[C]{\smallbreak\pagebreak[2]}
\correspDesc{Versand  durch Arthur Schnitzler, Olga Schnitzler am 26. 6. 1908 in Seis am Schlern
\newline{}Übermittlung  in Seis am Schlern
\newline{}Erhalt  durch Gustav Schwarzkopf im Zeitraum [27. 6. 1908 – 1. 7. 1908?] in Wien}\toendnotes[C]{\smallbreak}
\Standort{DLA, A:Schnitzler, HS.1985.1.1897.}
\physDesc{Bildpostkarte, 137 Zeichen
\newline{}Handschrift Arthur Schnitzler: schwarze Tinte, deutsche Kurrent
\newline{}Handschrift Olga Schnitzler: schwarze Tinte
\newline{}Versand: Stempel: »\nobreak{}\oindex{Seis am Schlern@\textbf{Seis am Schlern}|pwk}Seis\nobreak{}«.  }\pstart{}{\pb}Hrn \textsc{Gustav Schwarzkopf}\pend{}\pstart{}\textsc{Wien I}\oindex{I., Innere Stadt@\textbf{I., Innere Stadt}|pw}\pend{}\pstart{}\textsc{Tiefer Graben 23}\oindex{Wien@\textbf{Wien}!I., Innere Stadt@\textbf{I., Innere Stadt}!Tiefer Graben 23@\textbf{Tiefer Graben 23}|pw}.\pend{}{\bigskip}
\pstart
           \centering{}{\pb}\textcolor{gray}{\textbf{Tirol\oindex{Südtirol@\textbf{Südtirol}|pw}: Seis am Schlern\oindex{Seis am Schlern@\textbf{Seis am Schlern}|pw}, 1000 m. Nach dem Aquarell von F. A. C. M. Reisch\pwindex{Reisch, Franz August Carl Maria 1.\,5.\,1862 Wien – 1942? Meran@\textsc{Reisch, Franz August Carl Maria} (1.\,5.\,1862 Wien – 1942? Meran)|pw}, Meran\oindex{Meran@\textbf{Meran}|pw}.}}\pend
           
\pstart
           \centering{}\textcolor{gray}{\textbf{Santnerspitze}}\oindex{Punta Santner@\textbf{Punta Santner}|pw}\hspace*{2.5em}\textcolor{gray}{\textbf{Burgstall}}\oindex{Burgstall@\textbf{Burgstall}|pw}\hspace*{2.5em}\textcolor{gray}{\textbf{Jung-Schlern\oindex{Piccolo Sciliar@\textbf{Piccolo Sciliar}|pw}}}\pend
           \vspace{1em}
\pstart
           {\pb}Villa Heufler\oindex{Villa Heufler@\textbf{Villa Heufler}|pw},
                     26. 6. 08\pend
           \vspace{0.5em}
\pstart
           {\pb}Herzlichſte Grüße, es gefällt uns hier ausnehmend gut.\pend
           
\pstart
           Ihr \spacefill\mbox{Arthur}{\\[\baselineskip]}\spacefill\mbox{{[}hs. O. Schnitzler:{]} Olga.}\pend
           \leftskip=0em{}\selectlanguage{ngerman}\endnumbering\briefempfaengerindex{Schwarzkopf, Gustav@\textsc{Schwarzkopf, Gustav}!zzzSchnitzler, Olga@\emph{von Olga Schnitzler}!1908-06-262@{26. 6. 1908}|)be}\briefempfaengerindex{Schwarzkopf, Gustav@\textsc{Schwarzkopf, Gustav}!zzzSchnitzler, Arthur@\emph{von Arthur Schnitzler}!1908-06-262@{26. 6. 1908}|)be}\mylabel{L04373h}
\begin{anhang}
\end{anhang}\newcommand{\dateiname}{L04373}\newcommand{\titel}{Arthur und Olga Schnitzler an Gustav Schwarzkopf, 26. 6. 1908}\newcommand{\editorInnen}{Herausgegeben von Jahnke, SelmaMüller, Martin Anton}\input{../tex-inputs/latex-leseansicht-abspann}
